
\documentclass[11pt,a4paper]{article}

\usepackage[a4paper,margin=2.6cm]{geometry}
\usepackage{amsmath,amssymb,amsthm,mathtools}
\usepackage{booktabs}
\usepackage{enumitem}
\usepackage{xcolor}
\usepackage{microtype}
\usepackage[hidelinks]{hyperref}

\newcommand{\R}{\mathbb R}
\newcommand{\F}{\mathcal F}
\newcommand{\W}{\mathcal W}
\newcommand{\supp}{\operatorname{supp}}
\newcommand{\spread}{\operatorname{spread}}
\newcommand{\jointfloor}{\operatorname{jointfloor}}

\newcommand{\PROVED}{\textcolor{green!45!black}{\textbf{PROVED}}}
\newcommand{\NUM}{\textcolor{blue!70!black}{\textbf{NUMERICAL}}}
\newcommand{\OPEN}{\textcolor{red!70!black}{\textbf{OPEN}}}

\newtheorem{proposition}{Proposition}
\newtheorem{lemma}[proposition]{Lemma}
\newtheorem{remark}[proposition]{Remark}
\newtheorem{definition}[proposition]{Definition}

\title{\textbf{Raw Mass Is Not Routing Phase}\\
\large A normalization correction for antidiagonal optimal-face tomography}

\author{}
\date{30 August 2026}

\begin{document}
\maketitle

\begin{abstract}
A recent optimal-face audit of the quartic zero-repair linear program found a
strictly positive joint antidiagonal spread,
\[
\jointfloor^\star
=
1.475854331283591\times 10^{-2},
\]
together with an essentially exact dual certificate.  The certificate is
mathematically correct, but it reveals that the tested observable was the
\emph{raw antidiagonal routing mass}
\[
W_h(s,x)
=
\sum_{j+y=s} z_h(j,y,x),
\]
rather than the normalized conditional routing law
\[
G_h(x\mid s)
=
\frac{W_h(s,x)}{R_h(s)},
\qquad
R_h(s)=\sum_x W_h(s,x).
\]
In the reported instance the entire positive joint floor is explained by a
difference in total antidiagonal mass:
\[
\jointfloor^\star
=
\frac{1}{3}
\left|
R_{(3,0)}(2)-R_{(2,1)}(2)
\right|
\]
up to approximately \(3.5\times10^{-15}\).
Thus the computation rigorously detects unavoidable history dependence of
\emph{mass}, but it does not by itself prove unavoidable history dependence of
the \emph{normalized routing proportions}.  This note isolates the distinction,
proves the general pigeonhole certificate behind the observed floor, retracts
the stronger phase interpretation, and formulates the correct next tests.
\end{abstract}

\paragraph{Epistemic convention.}
Statements of finite-dimensional linear algebra are labelled \PROVED.
Reported solver values are labelled \NUM.
No finite numerical observation is promoted to an infinite-horizon theorem.

% ------------------------------------------------------------------
\section{Setup}
% ------------------------------------------------------------------

At an internal history \(h\), let
\[
z_h(j,y,x)
=
\Pr(J=j,Y=y,X=x)
\]
denote the zero-repair coupling variables.  Introduce the antidiagonal
coordinate
\[
s=j+y.
\]
For a fixed \(s\) and output symbol \(x\), define the raw routed mass
\begin{equation}
\boxed{
W_h(s,x)
:=
\sum_{j+y=s} z_h(j,y,x).
}
\tag{1}
\end{equation}
Summing over admissible \(x\) gives the total mass incident on the
antidiagonal:
\begin{equation}
\boxed{
R_h(s)
:=
\sum_x W_h(s,x)
=
\sum_{j+y=s} A_h(j)Q(y).
}
\tag{2}
\end{equation}
Whenever \(R_h(s)>0\), the normalized conditional routing law is
\begin{equation}
\boxed{
G_h(x\mid s)
:=
\frac{W_h(s,x)}{R_h(s)}.
}
\tag{3}
\end{equation}

Equations \((1)\)--\((3)\) describe two genuinely different notions of
history dependence:

\[
\boxed{
\begin{array}{ll}
\text{mass dependence:}
&
R_h(s)\ \text{depends on }h,
\\[1mm]
\text{routing-phase dependence:}
&
G_h(\cdot\mid s)\ \text{depends on }h.
\end{array}
}
\tag{4}
\]

The first may occur even when all histories use exactly the same normalized
routing proportions.

% ------------------------------------------------------------------
\section{What the joint-floor LP actually certifies}
% ------------------------------------------------------------------

Let \(\F^\star\) denote the chosen optimal or near-optimal face.  For a raw
cell \(q=(s,x)\), define
\[
\spread_q(z)
=
\max_h W_h(s,x)-\min_h W_h(s,x).
\]
The joint-face program computes
\begin{equation}
\boxed{
\jointfloor^\star
=
\min_{z\in\F^\star}
\max_q \spread_q(z).
}
\tag{5}
\end{equation}
This is a linear program after introducing upper and lower envelope variables
for each \(q\).

The important point is that \((5)\) is a statement about \(W\), not \(G\).
A positive value therefore proves that no point of \(\F^\star\) can make all
\emph{raw} routing masses history-independent.  It does not yet identify why.

The extracted dual certificate resolves that question in the tested instance.

\begin{proposition}[Total-mass pigeonhole bound]
\label{prop:pigeonhole}
Fix an antidiagonal \(s\) with \(m\) admissible output symbols
\(x\in X_s\), and two histories \(h_+\) and \(h_-\).  Then
\begin{equation}
\boxed{
\max_{x\in X_s}
\left|
W_{h_+}(s,x)-W_{h_-}(s,x)
\right|
\ge
\frac{
\left|R_{h_+}(s)-R_{h_-}(s)\right|
}{m}.
}
\tag{6}
\end{equation}
Consequently,
\begin{equation}
\boxed{
\max_{x\in X_s}
\spread_{s,x}(z)
\ge
\frac{
\left|R_{h_+}(s)-R_{h_-}(s)\right|
}{m}.
}
\tag{7}
\end{equation}
\end{proposition}

\begin{proof}
Write
\[
d_x
=
W_{h_+}(s,x)-W_{h_-}(s,x).
\]
By definition,
\[
\sum_{x\in X_s} d_x
=
R_{h_+}(s)-R_{h_-}(s).
\]
Hence
\[
\left|R_{h_+}(s)-R_{h_-}(s)\right|
=
\left|\sum_x d_x\right|
\le
\sum_x |d_x|
\le
m\max_x |d_x|,
\]
which proves \((6)\).  Since
\[
\spread_{s,x}(z)
\ge
\left|
W_{h_+}(s,x)-W_{h_-}(s,x)
\right|,
\]
equation \((7)\) follows.
\end{proof}

\paragraph{Status.}
Proposition~\ref{prop:pigeonhole} is \PROVED.

% ------------------------------------------------------------------
\section{The observed certificate}
% ------------------------------------------------------------------

In the depth-\(3\), \(M=8\) audit, the active dual support selected the three
cells
\[
(s,x)\in\{(2,0),(2,1),(2,2)\}
\]
with equal weights
\[
\boxed{
\alpha_0=\alpha_1=\alpha_2=\frac13,
}
\tag{8}
\]
and the history pair
\[
\boxed{
h_+=(3,0),
\qquad
h_-=(2,1).
}
\tag{9}
\]
Since \(x=0,1,2\) are exactly the admissible output symbols on \(s=2\), their
sum is the full antidiagonal mass:
\[
\sum_{x=0}^2 W_h(2,x)=R_h(2).
\]
The extracted certificate therefore reduces to
\begin{equation}
\boxed{
\jointfloor^\star
\ge
\frac13
\left|
R_{(3,0)}(2)-R_{(2,1)}(2)
\right|.
}
\tag{10}
\end{equation}

The reported numerical values are
\begin{align}
\jointfloor^\star
&=
0.01475854331283591,
\tag{11}\\
\frac13
\left|
R_{(3,0)}(2)-R_{(2,1)}(2)
\right|
&=
0.01475854331283239.
\tag{12}
\end{align}
Their difference is approximately
\[
\boxed{
3.5\times10^{-15}.
}
\tag{13}
\]
The KKT stationarity residual of the extracted dual certificate is of the
same order.

\paragraph{Status.}
Equations \((11)\)--\((13)\) are \NUM.  They show, to essentially machine
precision, that the measured joint floor is saturated by the total-mass
pigeonhole obstruction.

This is an important \emph{explanation} of the finite LP value, but it changes
the interpretation of that value.

% ------------------------------------------------------------------
\section{Correction of the phase interpretation}
% ------------------------------------------------------------------

The previously tempting interpretation was:

\begin{quote}
A positive joint floor proves that some history-dependent fresh routing phase
survives quotienting by optimal-face degeneracy.
\end{quote}

That statement is too strong for the observable actually tested.

Indeed, suppose for some fixed probability vector \(p_s(x)\) that
\begin{equation}
\boxed{
W_h(s,x)=R_h(s)\,p_s(x)
\qquad
\text{for every history }h.
}
\tag{14}
\end{equation}
Then
\[
G_h(x\mid s)=p_s(x)
\]
is completely history-independent, while the raw masses \(W_h(s,x)\) still
vary whenever \(R_h(s)\) varies.

Thus
\[
R_{h_+}(s)\neq R_{h_-}(s)
\]
can force a strictly positive raw joint floor even in a world with
\emph{zero normalized routing phase}.

The corrected conclusion is therefore
\begin{equation}
\boxed{
\jointfloor^\star(W)>0
\quad\Longrightarrow\quad
\text{unavoidable history dependence of raw antidiagonal mass,}
}
\tag{15}
\end{equation}
but not
\begin{equation}
\boxed{
\jointfloor^\star(W)>0
\quad\not\Longrightarrow\quad
\text{unavoidable history dependence of }G_h(\cdot\mid s).
}
\tag{16}
\end{equation}

\paragraph{Status.}
The implication distinction \((15)\)--\((16)\) is \PROVED.

% ------------------------------------------------------------------
\section{Why this does not make the computation useless}
% ------------------------------------------------------------------

The joint-floor computation remains mathematically valuable for three reasons.

First, it correctly exposed a nontrivial fact about the entire optimal face:
the raw antidiagonal mass profile cannot be made simultaneously
history-independent.

Second, the dual extractor did more than reproduce a number.  It identified
the exact finite-dimensional mechanism behind the floor: a three-channel
partition of one antidiagonal together with a total-mass imbalance between two
histories.

Third, the correction cleanly separates two levels of structure that should
not be conflated:
\[
\boxed{
A_h
\ \longrightarrow\
R_h(s)
\ \longrightarrow\
W_h(s,x)
=
R_h(s)G_h(x\mid s).
}
\tag{17}
\]
History dependence can enter already at the \(A_h\) or \(R_h\) level.
A claim of genuinely new routing phase must survive after that amplitude
dependence has been controlled.

% ------------------------------------------------------------------
\section{The correct next tests}
% ------------------------------------------------------------------

\subsection{A linear mass-controlled slice}

For a selected pair of histories \(h_+,h_-\) and antidiagonal \(s\), impose
the additional linear equality
\begin{equation}
\boxed{
R_{h_+}(s)=R_{h_-}(s).
}
\tag{18}
\end{equation}
If the resulting face slice is nonempty, then any remaining raw routing
difference between those two histories cannot be explained by a difference in
their total antidiagonal mass.

Define
\begin{equation}
\Phi^{\mathrm{eqmass}}_{h_+,h_-;s}
:=
\min_{\substack{z\in\F^\star\\
R_{h_+}(s)=R_{h_-}(s)}}
\max_{x\in X_s}
\left|
W_{h_+}(s,x)-W_{h_-}(s,x)
\right|.
\tag{19}
\end{equation}
After the usual epigraph linearization, \((19)\) is an LP.

If
\[
\Phi^{\mathrm{eqmass}}_{h_+,h_-;s}>0,
\]
then normalized routing differs between \(h_+\) and \(h_-\) on that
equal-mass slice, because their denominators are identical.

This is a direct and inexpensive falsification test.

\subsection{Direct normalized routing}

The intrinsic quantity is
\[
G_h(x\mid s)=\frac{W_h(s,x)}{R_h(s)}.
\]
Single-ratio extrema such as
\[
\min_{z\in\F^\star} G_h(x\mid s),
\qquad
\max_{z\in\F^\star} G_h(x\mid s)
\]
are linear-fractional programs and may be converted to LPs by a
Charnes--Cooper transformation when \(R_h(s)>0\) is controlled.

The genuinely history-comparative condition
\[
G_h(x\mid s)=G_g(x\mid s)
\]
is equivalent to
\[
W_h(s,x)R_g(s)=W_g(s,x)R_h(s),
\]
which is bilinear.  A complete optimal-face test of simultaneous equality of
normalized routing therefore requires either a dedicated bilinear/global
optimization formulation or a suitable exact reformulation.

\paragraph{Status.}
The equal-mass LP \((19)\) is a concrete \PROVED\ reduction.
A complete normalized-routing face tomography is \OPEN.

% ------------------------------------------------------------------
\section{Revised interpretation of the ``ghost'' cell}
% ------------------------------------------------------------------

The same normalization warning applies to the previously reported positive
floor at
\[
(j,y,x)=(6,1,4).
\]
A positive raw spread there proves that the corresponding coupling mass cannot
be made history-independent over the tested face.  It does not, by itself,
separate a change in total incident mass from a change in normalized routing
preference.

Accordingly, the phrase ``fresh causal phase is proved to exist'' should not
be used for that observation unless the effect survives a mass-controlled or
normalized-routing test.

The safe statement is:

\begin{quote}
The finite optimal face contains solver-independent history dependence in raw
coupling mass at the reported cell.  Whether this represents genuinely
history-dependent normalized routing remains unresolved.
\end{quote}

% ------------------------------------------------------------------
\section{Implication for the upper-bound program}
% ------------------------------------------------------------------

This correction does not settle the logarithmic upper bound in either
direction.

A history-dependent \(R_h(s)\) may be entirely compatible with a bounded or
logarithmic total bill.  Conversely, even a small normalized-routing phase
could in principle carry a large transport cost.  The upper-bound question
therefore cannot be decided from raw joint floors alone.

The useful consequence is methodological:

\begin{equation}
\boxed{
\text{Any future antidiagonal ``phase'' observable must quotient or control
the total-mass coordinate }R_h(s).
}
\tag{20}
\end{equation}

Only the residual structure after this normalization is a candidate for a
genuinely new causal routing obstruction.

% ------------------------------------------------------------------
\section{Research protocol}
% ------------------------------------------------------------------

For future finite-depth audits we propose the following order.

\begin{enumerate}[leftmargin=2em]
\item Compute the raw optimal-face observable \(W_h(s,x)\).
\item Compute and report \(R_h(s)=\sum_xW_h(s,x)\) on the same face.
\item Before interpreting a positive raw spread as routing phase, test the
      equal-mass slice \((18)\).
\item Where feasible, compute direct linear-fractional bounds on
      \(G_h(x\mid s)\).
\item Promote a ``fresh phase'' claim only if history dependence survives
      normalization by \(R_h(s)\).
\item Keep the raw-mass and normalized-routing ledgers separate.
\end{enumerate}

This protocol prevents amplitude variation from being misidentified as phase
variation.

% ------------------------------------------------------------------
\section{Bottom line}
% ------------------------------------------------------------------

The joint-floor dual certificate is not wrong.  It is, in fact, sufficiently
clean to diagnose the earlier overinterpretation.

What was proved numerically on the finite face is
\[
\boxed{
\text{unavoidable history dependence of raw antidiagonal mass.}
}
\]
What remains to be proved is
\[
\boxed{
\text{unavoidable history dependence of normalized antidiagonal routing.}
}
\]

The distinction is exactly
\[
\boxed{
W_h(s,x)
=
R_h(s)\,G_h(x\mid s).
}
\]

The next experiment should therefore not ask whether \(W\) differs across
histories.  It should ask whether any difference survives after the
\(R\)-coordinate has been removed.

\end{document}
